\subsection{李代数上的积分}

命题 2. 令 H 为一个（实）维数为 m 的李群，h 为其李代数。令 \(\omega_{H}\) 为 H 上的 m 次右不变微分形式，并令 \(\omega_{h}\) 为 h 上的 m 次平移不变微分形式，它在原点处与 \(\omega_{\mathrm{H}}(e)\) 相同。我们有

\[
(\exp_ {\mathrm{H}}) ^ {*} \omega_ {\mathrm{H}} = \lambda_ {\mathfrak {h}} \omega_ {\mathfrak {h}}\tag{10}
\]

其中 \(\lambda_{\mathfrak{h}}\) 是 \(\mathfrak{h}\) 上的 Ad(H)-不变函数，满足

\[
\lambda_ {\mathfrak {h}} (x) = \det \left(\sum_ {p \geq 0} \frac {1}{(p + 1) !} (\operatorname{ad} x) ^ {p}\right) \quad \text { for } x \in \mathfrak {h}.
\]

令 \(x, x_1, \ldots, x_m\) 为 \(\mathfrak{h}\) 中的元素。我们有

\[
(\exp^ {*} \omega_ {\mathrm{H}}) _ {x} (x _ {1}, \dots , x _ {m}) = (\omega_ {\mathrm{H}} (\exp x)) (\mathrm{T} _ {x} (\exp) (x _ {1}), \dots , \mathrm{T} _ {x} (\exp) (x _ {m})).
\]

将 \(\varpi(x) : \mathfrak{h} \to \mathfrak{h}\) 记为指数映射在 \(x\) 处的右微分（第 III 章，第 3 节，第 17 号，定义 8）；按定义，

\[
\mathrm{T} _ {x} (\exp) (y). (\exp x) ^ {- 1} = \varpi (x). y \quad \text { for   all } y \in \mathfrak {h}.
\]

由于形式 \(\omega_{H}\) 是右不变的，我们得到

\[
\begin{array}{r l} & (\omega_ {\mathrm{H}} (\exp x)) (\mathrm{T} _ {x} (\exp) (x _ {1}), \ldots , \mathrm{T} _ {x} (\exp) (x _ {m})) \\ & \quad = \omega_ {\mathrm{H}} (e) (\varpi (x). x _ {1}, \ldots , \varpi (x). x _ {m}) = (\det \varpi (x)) \omega_ {\mathfrak {h}} (x _ {1}, \ldots , x _ {m}); \end{array}
\]

因此，\(\exp^{*}\omega_{\mathrm{H}} = \lambda_{\mathfrak{h}}\omega_{\mathfrak{h}}\)，其中 \(\lambda_{\mathfrak{h}}(x) = \det \varpi(x) = \det \frac{\exp\operatorname{ad} x - 1}{\operatorname{ad} x}\)（第 III 章，第 6 节，第 4 号，命题 12）。

令 \(h \in H\)；由于 Ad h 是 h 的自同构，有

\[
\operatorname{ad} \left((\operatorname{Ad} h) (x)\right) = \operatorname{Ad} h \circ \operatorname{Ad} x \circ (\operatorname{Ad} h) ^ {- 1},
\]

所以 \(\lambda_{\mathfrak{h}}((\mathrm{Ad}h)(x)) = \lambda_{\mathfrak{h}}(x)\)。因此，函数 \(\lambda_{\mathfrak{h}}\) 在 \(\mathrm{Ad(H)}\) 下不变；命题得证。

注. 考虑与紧李群 G 关联的函数 \(\lambda_{g}\)；根据第 2 节第 1 号定理，要计算 \(\lambda_{g}\) 只需知道它在 t 上的值。但是，对于 \(x \in t\)，g 的自同态 ad x 是半单的，其特征值为 0（重数为 r），以及对于所有 \(\alpha \in \mathrm{R}(\mathrm{G}, \mathrm{T})\) 的 \(\delta(\alpha)(x)\)（重数为 1）。立即得到

\[
\lambda_ {\mathfrak {g}} (x) = \prod_ {\alpha \in \mathrm{R} (\mathrm{G}, \mathrm{T})} \frac {e ^ {\delta (\alpha) (x)} - 1}{\delta (\alpha) (x)} = \frac {\delta_ {\mathfrak {g}} (x)}{\pi_ {\mathfrak {g}} (x)}\tag{11}
\]

\[
\text { with } \delta_ {\mathfrak {g}} (x) = \delta_ {\mathrm{G}} (\exp x) \text { and } \pi_ {\mathfrak {g}} (x) = \prod_ {\alpha \in \mathrm{R} (\mathrm{G}, \mathrm{T})} \delta (\alpha) (x) = \det \operatorname{ad} _ {\mathfrak {g} / \mathfrak {t}} (x).
\]

令 \(\omega_{\mathrm{G / T}}\) 为次数 \(n - r\) 的 \(\mathrm{G / T}\) 上不变微分形式，令 \(\omega_{\mathfrak{t}}\) 为次数 \(r\) 的 \(\mathfrak{t}\) 上平移不变微分形式。按第 1 号的记号，将 \(\omega_{\mathrm{G / T}} \cap \omega_{\mathfrak{t}}\) 记为唯一的平移不变微分形式 \(\omega_{\mathfrak{g}}\)，其次数为 \(n\)，定义在 \(\mathfrak{g}\) 上，且满足 \(\omega_{\mathfrak{g}}(0) = \omega_{\mathrm{G / T}}(\bar{e}) \cap \omega_{\mathfrak{t}}(0)\)。

最后，将 \(\psi:(\mathrm{G}/\mathrm{T})\times\mathfrak{t}\to\mathfrak{g}\) 记为由映射 \((g,x)\mapsto(\operatorname{Ad}g)(x)\) 从 \(G\times t\) 到 g 通过取商诱导的流形态射。

命题 3. 设 \(\omega_{\mathfrak{g}}\)、\(\omega_{\mathfrak{t}}\)、\(\omega_{\mathrm{G / T}}\) 分别是 \(\mathfrak{g}\)、\(\mathfrak{t}\)、\(\mathrm{G / T}\) 上次数分别为 \(n, r, n - r\) 的不变微分形式。若 \(\omega_{\mathfrak{g}} = \omega_{\mathrm{G / T}} \cap \omega_{\mathfrak{t}}\)，则有

\[
\psi^ {*} \omega_ {\mathfrak {g}} = \omega_ {\mathrm{G/T}} \wedge \pi_ {\mathfrak {g}} \omega_ {\mathfrak {t}}
\]

其中 \(\pi_{\mathfrak{g}}\) 是 \(\mathfrak{t}\) 上由 \(\pi_{\mathfrak{g}}(x) = \prod_{\alpha \in \mathrm{R}(\mathrm{G},\mathrm{T})}\delta (\alpha)(x)\) 定义的函数。

将 \(\omega_{G}\)（分别为 \(\omega_{T}\)）记为 G（分别为 T）上最高次数的不变微分形式，它们在原点处分别与 \(\omega_{g}\)（分别为 \(\omega_{t}\)）相同。考虑交换图

\[
\begin{array}{c c c} (\mathrm{G} / \mathrm{T}) \times \mathfrak {t} & \stackrel {{\psi}} {{\longrightarrow}} & \mathfrak {g} \\ \Big \downarrow^ {(\mathrm{Id}, \exp_ {\mathrm{T}})} & & \Big \downarrow^ {\exp_ {\mathrm{G}}} . \\ (\mathrm{G} / \mathrm{T}) \times \mathrm{T} & \stackrel {{f}} {{\longrightarrow}} & \mathrm{G} \end{array}
\]

根据第 2 号的命题 1 以及关系 \(\exp_{T}^{*}\omega_{T} = \omega_{t}\)，得到等式

\[
\psi^ {*} \mathrm{exp} _ {\mathrm{G}} ^ {*} \omega_ {\mathrm{G}} = \omega_ {\mathrm{G/T}} \wedge \delta_ {\mathfrak {g}} \omega_ {\mathfrak {t}}.
\]

根据命题 2，\(\psi^{*}\exp_{\mathrm{G}}^{*}\omega_{\mathrm{G}} = (\psi^{*}\lambda_{\mathfrak{g}})\psi^{*}\omega_{\mathfrak{g}}\)。由于函数 \(\lambda_{\mathfrak{g}}\) 在 \(\operatorname{Ad}(\mathrm{G})\) 下不变，有

\[
(\psi^ {*} \lambda_ {\mathfrak {g}}) (\bar {g}, x) = \lambda_ {\mathfrak {g}} (x) = \frac {\delta_ {\mathfrak {g}} (x)}{\pi_ {\mathfrak {g}} (x)} \quad \text { for } \bar {g} \in \mathrm{G/T}, x \in \mathfrak {t}.
\]

由此可知，形式 \(\psi^{*}\omega_{\mathrm{G}}(\bar{g},x)\) 与 \(\omega_{\mathrm{G / T}}(\bar{g})\wedge \pi_{\mathfrak{g}}(x)\omega_{\mathfrak{t}}(x)\) 在 \(\delta_{\mathfrak{g}}(x)\) 非零处相同，也就是在稠密开子集 \((\mathrm{G / T})\times \mathfrak{t}_r\) 上相同；因此它们处处相等，命题得证。

在 G 和 T 上分别选取最高次数的不变微分形式 \(\omega_{G}\) 和 \(\omega_{T}\)，使得 \(|\omega_{G}| = dg\) 且 \(|\omega_{T}| = dt\)；将 \(\omega_{g}\)（分别为 \(\omega_{t}\)）记为 g（分别为 t）上在原点处与 \(\omega_{\mathrm{G}}(e)\)（分别为 \(\omega_{\mathrm{T}}(e)\)）相同的平移不变微分形式，并将 Haar 测度 \(|\omega_{g}|\)（分别为 \(|\omega_{t}|\)）记为 dz（分别为 dx）。作与第 2 号类似并作相应修改的推理，得到以下命题：

命题 4. 测度 \(dz\)（定义在 \(\mathfrak{g}\) 上）是映射 \((g,x)\mapsto (\operatorname {Ad}g)(x)\) 从 \(\mathrm{G}\times \mathfrak{t}\) 到 \(\mathfrak{g}\) 的真映射下测度 \(dg\otimes \frac{1}{w(\mathrm{G})}\pi_{\mathfrak{g}}dx\) 的像。

第 2 号的推论 1 至 3 以及备注 1 至 3 的类似命题的叙述和证明留给读者。例如，设 \(\varphi\) 是 \(\mathfrak{g}\) 上的可积函数（取值于 Banach 空间或 \(\overline{\mathbf{R}}\)）；则

\[
\int_ {\mathfrak {g}} \varphi (z) d z = \frac {1}{w (\mathrm{G})} \int_ {\mathfrak {t}} \left(\int_ {\mathrm{G}} \varphi ((\mathrm{Ad} g) x) d g\right) \pi_ {\mathfrak {g}} (x) d x,\tag{12}
\]

特别地，若 \(\varphi\) 在 \(\operatorname{Ad}(\mathbf{G})\) 下不变，则

\[
\int_ {\mathfrak {g}} \varphi (z) d z = \frac {1}{w (\mathrm{G})} \int_ {\mathfrak {t}} \varphi (x) \pi_ {\mathfrak {g}} (x) d x.\tag{13}
\]

